\subsection{Compensación hacia atrás y concentración}

Fijamos un tiempo entero $T\le T_\zeta$ y una raíz. Para el valor terminal
$\sum_eZ^\oc_{e,\varnothing}$ resolvemos hacia atrás, en tiempo continuo $t\in[0,T]$, la ecuación
$\partial_t\Phi+\cL_t\Phi=0$ en el espacio lineal generado por~\eqref{eq:patrones}, escribiendo
$\Phi_t=\sum_{e,J}\{c^\oc_{e,J}(t)Z^\oc_{e,J}+c^\di_{e,J}(t)Z^\di_{e,J}\}$. La solución tiene forma
cerrada. Con $\Delta y=(T-t)/N$ y $z=z_t$, sean
\[
 \theta_e=\frac{\pi_e\,\Delta y}{z^4},\qquad \xi_{e,i}=\frac{p_{d_i}\,\Delta y}{z};
\]
entonces
\begin{equation}\label{eq:compensador_directo}
 \Phi_t=\sum_{e\in\mathcal I_{rc}}\bigl(O_e+\theta_e\chi_e\bigr)
 \prod_{i=1}^4\bigl(V_{d_i}-\xi_{e,i}\chi_{d_i}\bigr),
\end{equation}
es decir, $c^\oc_{e,J}=\prod_{i\in J}(-\xi_{e,i})$ y $c^\di_{e,J}=\theta_e\prod_{i\in J}(-\xi_{e,i})$.
En efecto, en $t=T$ se anulan $\theta_e$ y los $\xi_{e,i}$, y queda el valor terminal. Como
$\partial_t\Delta y=\partial_tz=-1/N$, se tiene $\partial_t\xi_{e,i}=-\beta_i+\alpha\xi_{e,i}$ y
$\partial_t\theta_e=-k_e+4\alpha\theta_e$; derivando los productos,
\[
\begin{aligned}
 \partial_tc^\oc_{e,J}&=|J|\,\alpha\,c^\oc_{e,J}+\sum_{i\in J}\beta_i\,c^\oc_{e,J\setminus\{i\}},\\
 \partial_tc^\di_{e,J}&=(4+|J|)\,\alpha\,c^\di_{e,J}+\sum_{i\in J}\beta_i\,c^\di_{e,J\setminus\{i\}}
 -k_e\,c^\oc_{e,J},
\end{aligned}
\]
que son las ecuaciones hacia atrás: por~\eqref{eq:generador}, con ellas se anula el coeficiente de
cada patrón en $\partial_t\Phi+\cL\Phi$. Es la misma estructura del Lema~\ref{lem:compensador}. Como
$0\le\Delta y/z\le1$, $p_{d_i}\le1$, $\pi_e\le C/n$ y $z\ge\zeta$, se obtiene, uniformemente,
\begin{equation}\label{eq:coeficientes}
 |c^\oc|\le C_\zeta,\quad |\partial_tc^\oc|\le C_\zeta/n,\quad |c^\di|\le C_\zeta/n,\quad
 |\partial_tc^\di|\le C_\zeta/n^2,
\end{equation}
y además $|\partial_t^2c^\oc|=O_\zeta(n^{-2})$ y $|\partial_t^2c^\di|=O_\zeta(n^{-3})$. Los
coeficientes resuelven la ecuación en tiempo continuo y se evalúan en tiempos enteros; el error
de esta discretización se acota en la prueba siguiente.

\begin{lema}[Concentración de los absorbedores]\label{lem:conc_directa}
Para toda tolerancia fija $\gamma>0$, eligiendo $\eps$ suficientemente pequeña en el
Lema~\ref{lem:regularidad},
\begin{equation}\label{eq:conc_directa}
 \Pr\bigl\{|A_T^{rc}-\Phi_0|>\gamma n,\ \sigma'>T\bigr\}\le e^{-c_{\zeta,\gamma}n},
\end{equation}
donde $\sigma'=\inf\{t:E_t\ge\eps n\text{ o fallo}\}$ cumple $\Pr(\sigma'\le T_\zeta)\le
e^{-c_{\zeta,\eps}n}$. La unión sobre las $n^2$ raíces y los $O(n)$ tiempos conserva la cola.
\end{lema}

\begin{proof}
\emph{Deriva.} Sea $\Delta c=c(t+1)-c(t)$ para cada coeficiente. Por la ecuación hacia atrás,
$\sum c\,\cL Z=-\sum\partial_tc\,Z$, y por tanto
\[
 \E(\Phi_{t+1}-\Phi_t\mid\mathcal F_t)=\sum c\,\bigl[\E(\Delta Z\mid\mathcal F_t)-\cL Z\bigr]
 +\sum(\Delta c-\partial_tc)\,Z+\E\Bigl(\sum\Delta c\,\Delta Z\Bigm|\mathcal F_t\Bigr),
\]
con las sumas sobre los patrones de la raíz y los coeficientes evaluados en $t$. Como
$\sum_eO_e\le n$ y hay $n^2$ casillas, el Lema~\ref{lem:generador} y~\eqref{eq:coeficientes}
acotan el primer término por $C_\zeta(\eps+n^{-1})$. Los otros dos forman el error de
discretización: por las cotas de $\partial_t^2c$, el segundo es $O_\zeta(n^{-1})$; como
$\sum|\Delta Z^\oc|=O(1)$ y $\sum|\Delta Z^\di|=O(n)$ por raíz, el tercero también. Así
$|\E(\Phi_{t+1}-\Phi_t\mid\mathcal F_t)|\le C_\zeta(\eps+n^{-1})$.

\emph{Saltos.} Una diagonal de destino afecta a lo sumo a dos candidatas ocupadas, por unicidad
de filas y columnas. Cualquier vértice aparece en $O(n)$ patrones disponibles de esta raíz, con
coeficientes $O(1/n)$. Un nacimiento añade un término ocupado de coeficiente acotado. Junto con
el cambio de coeficientes, esto da $|\Delta\Phi|\le C_\zeta$, también en la transición de salida.

\emph{Martingala.} El Lema~\ref{lem:regularidad} con tolerancia $\eps/2$ acota
$\Pr(\sigma'\le T_\zeta)$. Congelamos
$\widetilde\Phi_s=\Phi_{s\wedge\sigma'}(\text{estado}_{s\wedge\sigma'})$,
incluido su argumento temporal. La cota de la deriva vale también en la transición de salida; la
martingala $\widetilde\Phi_s-\widetilde\Phi_0-\sum_{u<s}\E(\Delta\widetilde\Phi_u\mid\mathcal F_u)$
tiene saltos acotados, y el error acumulado de deriva es $C_\zeta(2n\eps+2)$. Se fija $\eps$ para
absorberlo en la tolerancia $\gamma$ antes de tomar $n$ grande, y Azuma da~\eqref{eq:conc_directa},
pues $\Phi_T=A_T^{rc}$.
\end{proof}

\begin{lema}[Valor inicial]\label{lem:valor_inicial}
Con $y=T/N$,
\begin{equation}\label{eq:valor_inicial}
 \Phi_0=a_n^{rc}(y):=y\sum_{e\in\mathcal I_{rc}}\pi_e\prod_{i=1}^4\bigl(1-y\,p_{d_i}\bigr).
\end{equation}
\end{lema}

\begin{proof}
Al inicio $O_e=0$ y $\chi_v=V_v=1$ en todas las líneas, y en $t=0$ se tiene $\Delta y=y$ y $z=1$;
por~\eqref{eq:compensador_directo}, $\Phi_0=\sum_e\theta_e\prod_i(1-\xi_{e,i})$, que
es~\eqref{eq:valor_inicial}. Lo mismo se obtiene por dualidad. La solución hacia delante del
sistema lineal ideal~\eqref{eq:generador} con esos datos es, en el tiempo $y$ y con $z=1-y$,
\[
 Z^{\di,*}_{e,J}=z^{4+|J|}\prod_{i\notin J}(1-yp_{d_i}),\qquad
 Z^{\oc,*}_{e,J}=y\,\pi_e\,z^{|J|}\prod_{i\notin J}(1-yp_{d_i}),
\]
como se comprueba derivando: la fuente $k_eZ^{\di,*}_{e,J}$ cancela $z^4$ y los demás términos
son exactamente los de~\eqref{eq:generador}. Sea $Z^*(t)$ el vector de estos valores con $y=t/N$,
y $c(t)$ el de los coeficientes de $\Phi_t$. Entonces $\frac{d}{dt}\langle c(t),Z^*(t)\rangle=0$,
porque $c$ resuelve la ecuación hacia atrás y $Z^*$ la ecuación hacia delante con el mismo
operador; así $\Phi_0=\langle c(0),Z^*(0)\rangle=\langle c(T),Z^*(T)\rangle
=\sum_eZ^{\oc,*}_{e,\varnothing}(T/N)$. Es la solución de un sistema lineal, no una
factorización de la ley verdadera.
\end{proof}

\subsection{Trayectoria límite}\label{subsec:trayectoria}

Como $x_*\to0$, los perfiles objetivo convergen al perfil lleno $\nu_1$, de vacancias
diagonales $u_D,u_A\ge\kappa$. En esta subsección el perfil objetivo lleva el subíndice $x_*$, y
los símbolos sin subíndice denotan el perfil lleno. Por el Teorema~\ref{thm:perfil}
y~\eqref{eq:balance}, $\pi_e=b_{x_*}v_e/n\,(1+O(n^{-1}))$, donde $b_{x_*}v_e$ es el valor de
$\nu_{x_*}$ en el centro de $e$, y $p_d=1-u_{d,x_*}+O(n^{-1})$, donde $u_{d,x_*}$ es la vacancia
de $d$ en el perfil objetivo; por el Teorema~\ref{thm:perfil}(c), $\nu_{x_*}\to\nu_1$ y
$u_{d,x_*}\to u_d$ uniformemente. Las cuadraturas son uniformes y las $O(n)$ exclusiones de $G$,
fila y columna aportan $O(1/n)$. Por tanto, uniformemente en la raíz,
\begin{equation}\label{eq:trayectoria}
 \frac{a_n^{rc}(y)}n\longrightarrow\tilde a(y;\bar r,\bar c):=
 y\int_{[0,1]^2}\nu_1(i,j)\prod_{d\in\mathcal D(r,c;i,j)}\bigl[1-y(1-u_d)\bigr]\,di\,dj .
\end{equation}
En el límite, $u_d$ es la vacancia del perfil lleno. Por~\eqref{eq:regularidad}, $|q_t-qz_t|\le\eps n$ salvo
probabilidad $e^{-cn}$, de modo que $1-y$ es, en probabilidad, el límite de la fracción
de piezas reales pendientes; no se sustituye el perfil objetivo por el perfil óptimo de masa
$y$. La
función $\tilde a$ cumple $\partial_y\tilde a=\int\nu_1\prod_d\vartheta_d
-y\int\nu_1\sum_d(1-u_d)\prod_{\ell\ne d}\vartheta_\ell$, con $\vartheta_d=1-y(1-u_d)$ y
$\tilde a(0)=0$ (nacimiento menos pérdida), y depende de la posición de la raíz. En el tablero
lleno,
\begin{equation}\label{eq:a_infinito}
 \tilde a(1;r,c)=\int\nu_1\prod_du_d=a(r,c)\ge\kappa^4>0,
\end{equation}
con $a$ como en la Sección~\ref{sec:intro}. Esto es un límite del perfil determinista: no exige
regularidad del proceso hasta su último paso.
